\chapter{}
Desmanchamos nossa vida em Teresina e voltamos para Araraquara.
Alugamos uma casa e nos preparamos para recomeçar do nada.
Vicente foi matriculado num colégio particular, onde, em breve, eu voltaria a lecionar, felizmente.
Não tínhamos mais o carro que tinha sido vendido para pagar dívidas com o banco e isso tornou complicados os primeiros dias, quando um monte de providências tinham que ser tomadas para restabelecer a rotina da casa.
João sugeriu que eu usasse um carro da Fábrica que permanecia ocioso na garagem a maior parte do tempo.
Usei o carro por duas ou três vezes até ser avisada por Geraldo, o gerente, de que não podia continuar a fazê-lo, pois, ao contrário do que João dissera, a tal viatura não pertencia à Fábrica e sim ao Reginaldo, a título de complementação do salário que ele se atribuíra por assumir os negócios da família.
Esse foi o início de uma série de desagradáveis surpresas, a que se sucederam outras tantas indelicadezas que acabaram por me convencer que, ao chegarmos de volta quebrados, como todos sabiam que estávamos, provocamos a impressão, em alguns, de que nos atiraríamos sobre a herança como carcarás sobre a presa.

Não quero pormenorizar aqui esse período terrível de mágoas e decepções recíprocas que nos dividiu e nos envergonhou a todos, os herdeiros de Mathias Filpi.
Sobretudo eu, que conhecia a todos muito bem e que  por isso mesmo podia antever com absoluta clareza o que aconteceria a partir da decisão desastrosa da minha mãe, não devia ter perdido o norte.
Em Teresina, com a ajuda do Irmão Pires, de muita leitura e de muito lidar com problemas de alunos, professores e famílias, eu tinha avançado muito na identificação dos meus processos inconscientes, das minhas fragilidades e conseguido, eu achava, superar muita coisa, crescendo na minha autoconfiança, no domínio de mim mesma.
Mas, ai de mim, era muito cedo para cantar vitória.
Fisgada pelo clima de irracionalidade que se estabeleceu nas discussões sobre a maldita partilha, eu me deixei levar, deixei-me atingir e revidei com toda a força das minhas frustrações acumuladas.
E só me dei conta disso muito tempo depois.

A situação que antepunha de um lado da rinha minha mãe e Reginaldo, entrincheirados atrás da parte, digamos, produtiva da herança, que eram a fábrica e o sítio e, de outro, o restante dos herdeiros, com a relativa neutralidade de Maria Lúcia, foi-se deteriorando rapidamente.
Meu irmão mais velho, treinado nas mesas de negociação da fase mais selvagem e predatória do capitalismo tupiniquim, encarava a coisa como um negócio, que ele e mamãe estavam convencidos de que ninguém mais tinha competência para conduzir.
João não aceitava tal juízo e nem o que ele entendia como rejeição materna.
Paulinho, como sempre foi do feitio dele, observava e esperava, aprendendo e armazenando munição.
Maria Lúcia, que Renatão, antes de morrer, tinha conseguido introduzir na fábrica como sócia minoritária, pisava em ovos para não perder essa condição que, ela acreditava, garantia-lhe a segurança de viúva recente.

As reuniões sucessivas para avaliar propostas e contrapropostas de ambos os lados para por fim às contendas, aconteciam em casa quase sempre, o que me mortificava, porque eu sabia que Paulo não tinha estrutura para aguentar aquele clima pesadíssimo.
As sucessivas referências à nossa situação, que pontificavam aqui e ali durante esses encontros, faziam-no sentir-se ainda pior, até porque, de fato, naquele instante, aquela herança significava a esperança de podermos nos reerguer.
Mas, duas outras agravantes contribuíam para desestabilizá-lo ainda mais: ver João e Paulinho, dos quais ele se sentia uma espécie de irmão mais velho, serem tratados como incapazes suscitava nele, acredito, terríveis lembranças da sua própria relação com o pai e, por outro lado, sentir minha resistência à ideia de pôr-lhe nas mãos os bens que me caberiam na partilha, identificava-me com os Filpi no conceito que faziam dele.

Nunca consegui que Paulo entendesse o real significado daquela herança para mim.
Acabei desistindo de tentar e lhe entreguei tudo.
Mas o que eu achava é que, se era minha necessidade de segurança o motivo mais sério para os nossos desentendimentos, com aquele dinheiro eu poderia garantir os estudos dos meninos, causa maior da minha ansiedade.
Além disso, libertado do compromisso de bancar os gastos da casa, ele poderia empreender o quisesse, sem cobranças, sem pressões e aí, quem sabe, suas iniciativas pudessem se consolidar.
A verdade é que essa seria a primeira oportunidade que eu teria, em toda minha vida, de tomar minhas próprias decisões acerca da minha vida e de garantir minha independência.
No fundo, eu precisava muito disso para me sentir inteira, forte, mas só descobriria a falta que isso me fazia quando, muitos anos depois, fui capaz de me sustentar sozinha, por quatro felizes anos.
A certeza, que por tanto tempo carreguei, de que nunca seria capaz fazê-lo, reduzia-me muito aos meus próprios olhos.

Paulo não entendeu, não aceitou minhas tentativas de esclarecer o que eu sentia, as brigas começaram a se suceder até que, num paroxismo, ele explodiu violentamente comigo, com o Reginaldo e, a partir daí, tudo se precipitou.
A partilha foi feita, tal como minha mãe e Reginaldo queriam, conseguimos evitar algumas injustiças, João e eu, não todas, mas a harmonia e a confiança que acreditávamos existir entre todos nós, que vivemos esse triste período, bateram asas, eu creio, para nunca mais voltar.
Os cacos de todas essas relações foram, ao longo dos anos que se seguiram, juntados aos poucos, colados, mas seriam para sempre como louça trincada.
Crescemos, mas as cicatrizes ficaram.

Paulo voltou a pensar na criação de peixes, dessa vez num projeto experimental em que eles seriam engordados em tanques de concreto, um sistema chamado de \textit{raceway}.
Arrendou uma área e começou a construir os tanques.
A implantação do projeto foi consumindo nossas reservas até que, para nosso alívio, surgiu-lhe uma oportunidade de trabalho em Assis, como gerente de compras de uma rede de supermercados.
O salário era bom, mas ele ficaria fora durante toda a semana, de modo que teve que deixar a condução do criatório a cargo do seu valete, Zezinho.

O entusiasmo recorrente do Paulo pela criação de peixes era e continua sendo um mistério para mim.
Porém, determinada a zelar pelo investimento dos nossos recursos, deixei minhas aulas e decidi supervisionar diariamente o trabalho do Zé.
Não que o empregado não fosse fiel ao Paulo ou desonesto, mas ele sempre me pareceu demasiadamente agitado para resignar-se àquele esforço de atenção concentrada.
Metida em macacão e botinas, para desespero da minha mãe, passei a frequentar a fazenda de segunda a sexta-feira, um período por dia, às vezes dois.
Nos primeiros tempos, até com uma das pernas metida em gesso porque, acompanhando Paulo numa viagem de treinamento para as novas funções, quebrei o pé.

\begin{figure}
\centering
\includegraphics[width=0.95\linewidth]{39/criando-tilápias-melhorada.jpeg}
\caption*{Criando tilápias em Araraquara.}
\end{figure}

\begin{figure}
\centering
\includegraphics[width=0.95\linewidth]{39/as-saint-peter-melhorada.jpeg}
\caption*{As ``Saint-Peter'' no tanque de engorda.}
\end{figure}

Trabalhei na fazenda de peixes durante cerca de um ano, cuidando dos berçários, ajudando na distribuição da ração nos tanques de engorda, na pesagem dos peixes para a venda.
Em nenhum momento me pareceu que se podia viver daquilo, ao menos não a curto ou médio prazo.
Os bichos eram frágeis, adoeciam à menor alteração de temperatura, a ração era cara e nem sempre eficiente e, naquele sistema criação intensiva, qualquer mínima contaminação podia ter consequências devastadoras.
Não foram poucas as vezes em que cheguei pela manhã e me deparei com um ou mais tanques em que os peixes boiavam moribundos, apesar de todos os nossos cuidados.
A fazenda onde o criatório estava instalado produzia laranjas e aplicava adubos e defensivos em grande escala.
Houve uma vez em que um capataz esqueceu aberta a comporta do rodolúvio em que os caminhões eram desinfetados e todo o líquido correu para os tanques.
Quando cheguei, pela manhã, não havia um único peixe vivo.


Muitas vezes Vicente, que na época cursava a faculdade de Engenharia na cidade vizinha, São Carlos, teve que acompanhar Zezinho ou o pai no meio da noite, ou na madrugada, ou ainda nos feriados, para acudir alguma ocorrência que pudesse comprometer a oxigenação dos tanques, nossa permanente preocupação.
Nessas ocasiões, ele costumava se lamentar:

\textit{``-- Por que eu não tenho um pai normalzinho, bancário, como tanta gente?''}

Apesar de tudo, não se pode negar, nossos peixes gozavam de alto conceito no mercado, porque os tanques de concreto lhes conferiam sabor delicado e o sistema de \textit{raceway} promovia uma engorda mais rápida, eficiente, e padronizada.
Nos dias em que tudo estava calmo, sentada ao lado dos tanques com um livro nas mãos, eu desfrutava preciosos momentos de serenidade, os primeiros em muitos anos.
Contudo, isso não me convencia de que nosso futuro pudesse depender daqueles serezinhos que me olhavam com aquele olhar arregalado e vazio.

Paulo, entretanto, seguia entusiasmado com a atividade e disposto até a ampliá-la.
A oportunidade surgiu quando um sócio capitalista se propôs a financiar essa expansão.
Aceita a sociedade, Paulo procurou uma nova área, muito maior, e iniciou a construção dos novos tanques.
Ele continuaria, ainda por uns tempos, a trabalhar em Assis.
Assim que a nova fazenda foi posta a funcionar, pediu demissão.
Logo depois, a ideia de vender peixes vivos foi ressuscitada.
Montamos dezoito grandes e belos aquários desenhados por Paulinho, meu irmão, e os distribuímos nas lojas da rede Carrefour, em São Paulo.
Os peixes eram pescados e abatidos diante dos clientes que podiam levá-los para casa com a certeza de consumir um produto absolutamente fresco e saudável.
Adquirimos um caminhão especial para transporte de peixes vivos que saía da fazenda noite ainda, para chegar às portas das lojas antes que elas se abrissem.
E logo, dezenas de pessoas, sobretudo de origem oriental, postavam-se à espera, ansiosas, diante dos aquários.
Um sucesso de vendas que chegou, nas palavras de alguns gerentes do setor, a elevar significativamente o desempenho das seções de pescados da rede francesa.
O sócio de Paulo estava animadíssimo.


\begin{figure}
\centering
\includegraphics[width=0.9\linewidth]{39/aquários-melhor.jpeg}
\caption*{Um dos nossos lindos aquários.}
\end{figure}

Uma madrugada, o telefone soou em nosso quarto.
Acontecia com frequência das lojas telefonarem pedindo mais peixes, com urgência.
Mas, daquela vez, o recado foi seco:

\textit{``-- Sr. Paulo, estamos tirando seus aquários da área de venda de todas as nossas lojas.
Por favor, providencie a retirada deles dos nossos depósitos o mais rapidamente possível.''}

Algumas senhoras com preocupações ecológicas e defensoras dos direitos dos animais, haviam se encarniçado contra o abate dos peixes diante do público e ameaçado ir aos jornais e à televisão contra a rede Carrefour.
Apavorados com a possibilidade, os diretores preferiram por fim imediato à iniciativa.

Desorientados, ensaiamos algum tipo de reação, mas não havia nada a fazer.
Só nos restava recolher e arquivar nossos belos aquários.

Paulo começou a procurar uma saída.
Na verdade não apenas para nós, porque os criadores de peixe, de modo geral, precisavam encontrar uma saída para a comercialização dos seus produtos que não fossem os pesque-e-pague que remuneravam mal os riscos e os custos da atividade.
Foi aí que Paulo pensou numa cooperativa que assumisse o processamento e a comercialização dos peixes, levando-os ao mercado sob a forma de filés, de hambúrgueres, de ``nuggets'', garantindo, enfim, seu total aproveitamento.

Os meses que se seguiram foram de uma atividade febril.
Paulo se empolgou, encontrou instalações propícias numa cooperativa de produtores de soja em Orlândia, negociou o arrendamento, foi atrás de equipamentos, realizou um monte de reuniões com os produtores, foi conhecer unidades de processamento no estado todo, providenciou o registro do frigorífico nos serviços de fiscalização, treinou funcionários e, finalmente, quase um ano depois, a cooperativa começou a funcionar.
Nesse tempo todo, pude acompanhá-lo muitas vezes.
Ele sempre se queixara de que eu nunca o apoiara e sempre o criticara nas suas empreitadas.
Dessa vez, desempregada e com os filhos crescidos, eu tinha tempo de sobra e decidi participar até onde me fosse possível.
E mais uma vez constatei o que eu já sabia.
Paulo ergueu a cooperativa com todos os requisitos para servir aos cooperados a melhor qualidade em serviços.
Como sempre, demonstrara a habitual facilidade em realizar uma obra, por mais complexa que pudesse parecer para a maioria das pessoas.
Mas, de novo, tropeçou no fator humano.
Tudo aquilo era fruto do que ele enxergava, lá na frente.
Mas os cooperados não o acompanharam.
Na realidade, nem chegavam a entender exatamente o espírito de uma cooperativa.
Não conseguiram ir adiante dos interesses particulares de cada um e acabaram por se sentir, é a minha impressão, arrastados pelo Paulo a instâncias que não sabiam se estavam dispostos a chegar e muito menos a bancar.
Sua liderança começou a ser contestada e ele acabou por entregar a presidência, retirando-se da cooperativa.
Muito pouco tempo depois, a cooperativa cessou suas atividades.
Concomitantemente, o sócio também pediu a dissolução da sociedade nas fazendas de peixes.
Ele vinha sendo mal sucedido numa aventura com criação de avestruzes e a empresa de fornecimento de mão-de-obra para a construção civil que lhe construíra a fortuna andava igualmente em crise.
Diante de tudo isso, não tendo como continuar, Paulo concordou com a liquidação do negócio.
A maior parte das instalações e equipamentos foi vendida para que se procedesse aos acertos com empregados, e fornecedores, e pagamento das contas.

Foi um duro golpe e eu me perguntava o que ele faria a seguir.
Uma amiga nossa, que montara um negócio de certificação de bois, contratou-o para ajudá-la a implantar todo o sistema de controles das fazendas produtoras.
O salário era bom e durante algum tempo ele se distraiu viajando e ajudando-a no projeto.
Mas, a parceria não iria longe.
Embora Paulo admirasse Cristina, essa nossa amiga, como empresária, ela era rude, e destemperada, e o próprio negócio a punha demasiadamente tensa.
Pedir paciência e compreensão ao Paulo, naquele momento, era exigir demais.
Ele preferiu, terminado o trabalho, desligar-se da firma e conservar a amizade.

